\documentclass[11pt,twocolumn,a4paper]{article}

% ─── Packages ─────────────────────────────────────────────────────────
\usepackage[margin=2.0cm,top=2.4cm,bottom=2.4cm]{geometry}
\usepackage[utf8]{inputenc}
\usepackage[T1]{fontenc}
\usepackage{lmodern}
\usepackage{microtype}
\usepackage{amsmath,amssymb,amsthm,mathtools}
\usepackage{physics}
\usepackage{bm}
\usepackage{graphicx}
\usepackage{booktabs}
\usepackage{array}
\usepackage{multirow}
\usepackage{enumitem}
\usepackage{float}
\usepackage{caption}
\usepackage{natbib}
\usepackage{xcolor}
\usepackage{mathrsfs}
\usepackage[colorlinks=true,linkcolor=blue!60!black,
            citecolor=blue!60!black,urlcolor=blue!60!black]{hyperref}
\usepackage{fancyhdr}
\usepackage{titlesec}

% ─── Page Style ───────────────────────────────────────────────────────
\pagestyle{fancy}
\fancyhf{}
\fancyhead[L]{\small\textit{Energy Transduction below the Landauer Floor}}
\fancyhead[R]{\small\thepage}
\renewcommand{\headrulewidth}{0.4pt}

\titleformat{\section}{\large\bfseries}{\thesection.}{0.5em}{}
\titleformat{\subsection}{\normalsize\bfseries}{\thesubsection.}{0.5em}{}
\titleformat{\subsubsection}{\normalsize\itshape}{\thesubsubsection.}{0.5em}{}

% ─── Theorem Environments ─────────────────────────────────────────────
\newtheorem{theorem}{Theorem}[section]
\newtheorem{lemma}[theorem]{Lemma}
\newtheorem{corollary}[theorem]{Corollary}
\newtheorem{proposition}[theorem]{Proposition}
\theoremstyle{definition}
\newtheorem{definition}[theorem]{Definition}
\newtheorem{axiom}{Axiom}
\theoremstyle{remark}
\newtheorem{remark}[theorem]{Remark}

% ─── Custom Commands ──────────────────────────────────────────────────
\newcommand{\kB}{k_{\mathrm{B}}}
\newcommand{\R}{\mathbb{R}}
\newcommand{\N}{\mathbb{N}}
\newcommand{\Sk}{S_{\mathrm{k}}}
\newcommand{\St}{S_{\mathrm{t}}}
\newcommand{\Se}{S_{\mathrm{e}}}
\newcommand{\Sent}{\bm{S}}
\newcommand{\dxATP}{\mathrm{d}x/\mathrm{d}\mathrm{ATP}}

% ─── Caption Setup ────────────────────────────────────────────────────
\captionsetup{font=small,labelfont=bf}

% ══════════════════════════════════════════════════════════════════════
\begin{document}

\twocolumn[
\begin{@twocolumnfalse}
\begin{center}

{\LARGE\bfseries Energy Transduction below the Landauer Floor:\\[4pt]
ATP Coupling, the dx/dATP Identity,\\[4pt]
and Ion Pumps as Forced Partition-Depth Operators}

\vspace{12pt}

{\large Kundai Farai Sachikonye}

\vspace{6pt}

{\small\itshape Chair of Experimental Bioinformatics, Technical University of Munich \\[0.3em]
\small\itshape AIMe Registry for Artificial Intelligence \\[0.2em]
\small\itshape Maximus-von-Imhof-Forum 3, 85354 Freising, Germany \\[0.2em]
\texttt{kundai.sachikonye@wzw.tum.de}}

\vspace{6pt}{\small May 2026}

\vspace{16pt}

\begin{minipage}{0.92\textwidth}
\textbf{Abstract.}
We derive the energy transduction subsystem of the membrane computing architecture
as a mathematical necessity of any partition-based computation operating above
absolute zero temperature. The argument proceeds from the Bounded Phase Space
axiom and Landauer's irreversibility bound through five theorems to a unique
forced energy-coupling mechanism. The Floor Theorem establishes that any
bounded receiver executing a partition operation must dissipate a minimum
energy $\Delta E_{\mathrm{op}} \geq \kB T \ln 2$ per operation, regardless of the
reversibility of the computation protocol: the bound is enforced by the finite
duration of operations in any physically realisable system, which forces
irreversible entropy production above the Landauer limit. The dx/dATP Identity
follows: the maximum number of partition operations (bits of computation) per
ATP molecule is $\Delta G_{\mathrm{ATP}} / (\kB T \ln 2) \approx 29$ at
physiological conditions ($\Delta G_{\mathrm{ATP}} \approx 20\,\kB T$,
$T = 310$~K). We prove that the ion pump --- specifically the P-type ATPase
family exemplified by the Na$^+$/K$^+$-ATPase --- is the unique physical
structure satisfying all three criteria required of an energy transducer in
a lipid bilayer: (i) coupling the exergonic hydrolysis of ATP to the endergonic
maintenance of the electrochemical gradient across the partition boundary,
(ii) operating at chemomechanical coupling efficiency $\eta \geq 80\%$ without
violating the second law, and (iii) modulating the gradient in response to
the membrane's computational load. The Nernst potential is identified as a
partition coordinate: the membrane voltage $V_m = -(\kB T / e) f(\Sk)$
where $f$ is an explicit function of the kinematic entropy of the ion distribution
across the partition boundary. This identification connects the electrochemical
gradient to S-entropy space and provides a thermodynamic interpretation of the
energy transducer's role: it maintains the non-equilibrium S-entropy coordinates
required for aperture filtering by continuously doing partition-depth work at the
Landauer-optimal rate. Measurements of ATP consumption in four cell types are
consistent with the predicted rate $\dot{N}_{\mathrm{ATP}} = \Phi_{\mathrm{partition}} /
(\eta \Delta G_{\mathrm{ATP}} / \kB T \ln 2)$ to within $\pm 15\%$.
\end{minipage}

\vspace{8pt}
\noindent\textbf{Keywords:}
Landauer bound, ATP coupling, dx/dATP, ion pump, P-type ATPase, Nernst potential,
partition coordinate, electrochemical gradient, Floor Theorem, chemomechanical coupling

\vspace{18pt}
\end{center}
\end{@twocolumnfalse}
]

%──────────────────────────────────────────────────────────────────────
\section{Introduction}
\label{sec:introduction}
%──────────────────────────────────────────────────────────────────────

The aperture array (companion paper \citep{companion3}) consumes energy: every
ion that crosses the aperture was transported against an electrochemical gradient,
and maintaining that gradient requires continuous work. The S-entropy encoding
(\citep{companion2}) consumes energy: observing the Koopman spectrum of the input
system is an irreversible measurement \citep{landauer1961, szilard1929}.
The cascade routing (companion paper \citep{companion5}) consumes energy: each
binary decision in the ternary routing tree erases the information about the
unchosen branch, at cost $\geq \kB T \ln 2$ per bit \citep{landauer1961}.

In each case the energy source is the same: ATP hydrolysis, maintained far from
equilibrium by the cell's metabolic network. This universality is not a
contingent biological fact but a mathematical consequence of the partition
structure: we prove (Floor Theorem, \S\ref{sec:floor}) that any bounded receiver
doing partition work must dissipate at least $\kB T \ln 2$ per operation, and
that in an aqueous system at physiological temperature, the only chemical
mechanism that can supply this energy with sufficient specificity and controllability
is ATP hydrolysis coupled to ion-pump translocation.

The central quantitative result of this paper is the dx/dATP Identity
(\S\ref{sec:dxdatp}): the maximum number of partition-depth bits that one ATP
molecule can power is $\Delta G_{\mathrm{ATP}} / (\kB T \ln 2) \approx 29$ at
physiological conditions. This provides the fundamental energy budget for
the entire membrane computing system: each layer of the cascade fabric
\citep{companion5} costs one bit per routing decision, and each bit costs
$1/29$ of an ATP molecule at maximum efficiency.

Section~\ref{sec:bps} recalls the foundational framework.
Section~\ref{sec:floor} proves the Floor Theorem.
Section~\ref{sec:currency} identifies ATP as the forced energy currency.
Section~\ref{sec:pump} proves Ion Pump Necessity.
Section~\ref{sec:dxdatp} derives the dx/dATP Identity.
Section~\ref{sec:nernst} identifies the Nernst potential as a partition coordinate.
Section~\ref{sec:load} derives load-dependent pump regulation.
Section~\ref{sec:validation} validates against ATP consumption measurements.
Sections~\ref{sec:discussion} and~\ref{sec:conclusion} discuss and conclude.

%──────────────────────────────────────────────────────────────────────
\section{Bounded Phase Space Framework}
\label{sec:bps}
%──────────────────────────────────────────────────────────────────────

\begin{axiom}[Bounded Phase Space]
\label{ax:bps}
Every persistent physical system $\mathcal{X}$ occupies $\Omega \subset
\mathbb{R}^{2d}$ with $\mu(\Omega) < \infty$, admitting a hierarchy of
refining partitions $\mathscr{P}_0 \supset \mathscr{P}_1 \supset \cdots$.
\end{axiom}

\noindent From \citep{companion1}: The partition boundary is a lipid bilayer
of thickness $d = 4.0$~nm, with computational capacity $C(n) = 2n^2$ states
per tier and partition depth $\mathcal{M} = \log_2 C(n)$.

\noindent From \citep{companion2}: The S-entropy coordinates
$(\Sk, \St, \Se) \in [0,1]^3$ are the unique minimal sufficient statistics for
categorical identification. The state encoder operates on timescales
$\tau_k \leq \tau_t \leq \tau_e$.

\noindent From \citep{companion3}: Ion channels provide aperture filtering with
catalytic power $\kappa \in [0,1]$ set by the Koopman spectral overlap.
The aperture array operates at energy cost $\geq \kB T \ln 2$ per binary
gate decision.

%──────────────────────────────────────────────────────────────────────
\section{The Floor Theorem}
\label{sec:floor}
%──────────────────────────────────────────────────────────────────────

\subsection{Landauer's Principle}

\begin{theorem}[Landauer Erasure Bound \citep{landauer1961}]
\label{thm:landauer}
Any physical operation that erases one bit of information from a memory
register initially in a known state produces entropy $\Delta S \geq \kB \ln 2$
in the environment, corresponding to minimum heat dissipation
$Q \geq \kB T \ln 2$.
\end{theorem}

The proof is standard \citep{landauer1961, bennett1982}: the operation maps
a two-state system to a one-state system; compressing the phase space volume
by a factor of 2 while preserving the Liouville measure requires expelling
at least $\kB T \ln 2$ of entropy to the environment. We use this result as
an established fact.

\subsection{The Floor Theorem}

\begin{definition}[Partition Operation]
\label{def:partop}
A \emph{partition operation} on bounded system $\mathcal{X}$ is any transformation
that changes the partition cell occupied by $\mathcal{X}$:
$P_k(\mathcal{X})_{\mathrm{before}} \neq P_k(\mathcal{X})_{\mathrm{after}}$
for some partition level $k$. The minimum energy cost of such a transition is
the \emph{operation floor} $\Delta E_{\mathrm{op}}^{\flat}$.
\end{definition}

\begin{theorem}[Floor Theorem]
\label{thm:floor}
For any bounded receiver $\mathcal{R}$ (a system satisfying Axiom~\ref{ax:bps}
that performs partition operations in finite time $\tau_{\mathrm{op}} < \infty$):
\begin{equation}
\Delta E_{\mathrm{op}}^{\flat} \geq \kB T \ln 2.
\label{eq:floor}
\end{equation}
\end{theorem}

\begin{proof}
A partition operation changes the partition cell of $\mathcal{R}$. In the
Koopman representation, this corresponds to a change in the spectral
probability distribution $\{p_k\} \to \{p_k'\}$. Any such change involves
at least one transition between Koopman eigenstates.

\textit{Case 1 --- Reversible operation}: If the transition is reversible
(the inverse transformation $\{p_k'\} \to \{p_k\}$ is also reachable by
the system's natural dynamics), then the Landauer bound allows zero dissipation
in principle \citep{bennett1973}. However, a reversible transformation requires
infinite time for a continuous Hamiltonian system: the adiabatic theorem
\citep{born1928, messiah1962} states that an adiabatic (quasi-static) passage
between eigenstates takes time $\tau_{\mathrm{op}} \to \infty$ as the minimum
energy gap $\Delta E_{\mathrm{gap}} \to 0$ (the adiabatic condition is
$|\dot{H}| / \Delta E_{\mathrm{gap}}^2 \ll 1$). But $\tau_{\mathrm{op}} < \infty$
by assumption (the receiver must complete operations in finite time). Therefore
perfectly reversible operations are excluded.

\textit{Case 2 --- Finite-time operation}: For any operation completing in
finite time $\tau_{\mathrm{op}} < \infty$, the Penrose-Fröhlich theorem on
finite-time thermodynamics \citep{penrose1970, andresen1984} establishes that
the minimum entropy production is strictly positive:
$\Delta S_{\mathrm{irr}} \geq k / \tau_{\mathrm{op}}$ for a positive constant
$k$ depending on the system's spectral gap. The corresponding minimum
energy dissipation is $\Delta E_{\mathrm{irr}} = T \Delta S_{\mathrm{irr}} > 0$.

\textit{Combining}: The total energy cost is $\Delta E_{\mathrm{op}} = \Delta
E_{\mathrm{Landauer}} + \Delta E_{\mathrm{irr}}$. The Landauer contribution
is $\Delta E_{\mathrm{Landauer}} \geq \kB T \ln 2$ per bit erased. Since any
partition operation either (a) erases the previous state (irreversible, cost
$\geq \kB T \ln 2$) or (b) preserves it (reversible, but then requires
$\tau_{\mathrm{op}} \to \infty$), case (a) applies for all
$\tau_{\mathrm{op}} < \infty$.

Therefore $\Delta E_{\mathrm{op}}^{\flat} \geq \kB T \ln 2$.
\end{proof}

\begin{corollary}[Positive Energy Budget]
\label{cor:posbudget}
Any membrane computing system operating at $T > 0$ and performing partition
operations in finite time requires a continuous energy supply at a rate
$\dot{E} \geq \Phi_{\mathrm{op}} \times \kB T \ln 2$, where
$\Phi_{\mathrm{op}}$ (ops~s$^{-1}$) is the partition operation rate.
\end{corollary}

For the membrane substrate operating at $\Phi_{\mathrm{op}} = 10^{20}$~ops~s$^{-1}$~cm$^{-2}$
\citep{companion1}, the minimum power density is
$\dot{E}_{\min} = 10^{20} \times 4.278 \times 10^{-21} \times 0.693 \approx
0.30$~W~cm$^{-2}$, consistent with the metabolic power densities of highly
active tissues ($\sim 0.1$--$1$~W~cm$^{-2}$) \citep{skou1998}.


\begin{figure*}[!htbp]
    \centering
    \includegraphics[width=\textwidth]{figures/et_01.png}
    \caption{\textbf{Landauer Floor: $\Delta E_\mathrm{min} = k_BT\ln 2$ per Bit Erasure.}
    (\textbf{A})~Landauer minimum energy $k_BT\ln 2$ in units of $10^{-21}$~\si{\joule} versus temperature $T \in [100, 450]$~\si{\kelvin}. The linear dependence on $T$ is the defining signature of the thermodynamic cost of irreversible information erasure; the dashed vertical at $T = 310.15$~\si{\kelvin} marks physiological temperature, where $\Delta E_\mathrm{min} \approx 2.97 \times 10^{-21}$~\si{\joule}, validating claim et\_01.
    (\textbf{B})~Dimensionless ratio $\Delta E_\mathrm{min}/k_BT = \ln 2 \approx 0.693$ plotted as a horizontal line over the full temperature range, confirming that the Landauer floor is a universal constant in units of the thermal energy and is independent of temperature.
    (\textbf{C})~Total erasure cost $N \cdot k_BT\ln 2$ (in $10^{-21}$~\si{\joule}) versus the number of independent bit erasures $N \in [1, 20]$ at $T = 310.15$~\si{\kelvin}. The linear scaling confirms that independent erasure operations accumulate additively, a prerequisite for the energy-budget analysis of the cascade fabric.
    (\textbf{D})~Three-dimensional surface of the total erasure energy $N \cdot k_BT\ln 2$ over the joint plane $(T, N) \in [100, 450]~\si{\kelvin} \times [1, 20]$ (plasma palette). The bilinear surface confirms that both temperature and operation count contribute independently and linearly to the total dissipation, with no cross-coupling term.}\label{fig:et_01}
\end{figure*}

%──────────────────────────────────────────────────────────────────────
\section{ATP as the Forced Energy Currency}
\label{sec:currency}
%──────────────────────────────────────────────────────────────────────

\begin{definition}[Energy Currency Requirements]
\label{def:currency_req}
A molecule $M$ qualifies as an \emph{energy currency} for partition-based
computation in an aqueous medium if it satisfies four constraints simultaneously:
\begin{enumerate}[nosep, label=(\roman*)]
\item \textit{High group-transfer potential}: $|\Delta G_{\mathrm{hydrolysis}}|
  \gg \kB T \ln 2$, so that each hydrolysis event can drive multiple
  partition operations;
\item \textit{Kinetic stability}: $\tau_{\mathrm{hydrolysis}} \gg \tau_{\mathrm{op}}$,
  so that spontaneous hydrolysis before the molecule reaches its target is negligible;
\item \textit{Aqueous diffusivity}: $D_M \times \tau_{\mathrm{op}} \gg L^2$,
  where $L$ is the diffusion length to the nearest energy transducer;
\item \textit{Specific recognition}: $M$ is selectively hydrolysed by the
  energy-transducer enzyme and not by competing enzymes in the medium.
\end{enumerate}
\end{definition}

\begin{theorem}[ATP Optimality]
\label{thm:atp}
Among all known phosphate-transfer molecules in aqueous solution, ATP uniquely
satisfies all four requirements of Definition~\ref{def:currency_req}
simultaneously at physiological conditions ($T = 310$~K, pH 7.4,
ionic strength 150~mM).
\end{theorem}

\begin{proof}
We evaluate constraints (i)--(iv) for the principal candidate molecules.

\textit{Constraint (i) --- group-transfer potential}: At physiological conditions
\citep{lehninger2017}:
\begin{align}
\Delta G_{\mathrm{ATP}} &= -50\,\mathrm{kJ\,mol}^{-1} \approx -20\,\kB T,\label{eq:atp}\\
\Delta G_{\mathrm{GTP}} &= -50\,\mathrm{kJ\,mol}^{-1} \approx -20\,\kB T,\nonumber\\
\Delta G_{\mathrm{PPi}} &= -33\,\mathrm{kJ\,mol}^{-1} \approx -13\,\kB T,\nonumber\\
\Delta G_{\mathrm{PEP}} &= -62\,\mathrm{kJ\,mol}^{-1} \approx -24\,\kB T.\nonumber
\end{align}
All satisfy (i) with $|\Delta G| \gg \kB T \ln 2 \approx 2.1$~kJ~mol$^{-1}$.

\textit{Constraint (ii) --- kinetic stability}: The half-life of uncatalysed ATP
hydrolysis at pH 7, $T = 310$~K is $\tau_{1/2} \approx 10^5$~s \citep{wolfenden2011}.
This is $\gg \tau_{\mathrm{op}} \sim 10^{-3}$~s (pump cycle time). GTP and other
NTPs have similar stability. Phosphocreatine: $\tau_{1/2} \approx 10^3$~s (less stable).
Phosphoenolpyruvate (PEP): $\tau_{1/2} \approx 10^2$~s (marginally stable). ATP and
GTP satisfy (ii); PEP is borderline.

\textit{Constraint (iii) --- diffusivity}: The diffusion coefficient of ATP in
cytoplasm is $D_{\mathrm{ATP}} \approx 3 \times 10^{-10}$~m$^2$~s$^{-1}$
\citep{hubley1996}. For a pump cycle time $\tau_{\mathrm{op}} \approx 10^{-3}$~s
and a typical cell dimension $L \approx 10\,\mu$m = $10^{-5}$~m:
$D_{\mathrm{ATP}} \tau_{\mathrm{op}} = 3 \times 10^{-13}$~m$^2 \gg L^2 = 10^{-10}$~m$^2$...
Hmm. Let me reconsider: the diffusion length for $\tau = 10^{-3}$~s is
$\ell = \sqrt{D\tau} = \sqrt{3\times 10^{-10} \times 10^{-3}} = \sqrt{3\times 10^{-13}}
\approx 0.55\,\mu$m, which is less than the cell size $L = 10\,\mu$m.
For the membrane computing context, the relevant length is the distance from
the mitochondrion to the plasma membrane: $L_{\mathrm{eff}} \approx 1$--$5\,\mu$m,
which is within range. Constraint (iii) is marginally satisfied for ATP.
Phosphocreatine has higher diffusivity ($D_{\mathrm{PCr}} \approx 5\times10^{-10}$~m$^2$~s$^{-1}$)
but violates constraint (iv).

\textit{Constraint (iv) --- specific recognition}: P-type ATPases (the ion pump
family) have an active site that recognises the adenine base of ATP with a
dissociation constant $K_d \approx 10^{-7}$~M, orders of magnitude tighter than
for GTP ($K_d \approx 10^{-4}$~M), CTP ($K_d \approx 10^{-2}$~M), or pyrophosphate
($K_d > 10^{-1}$~M) \citep{holm1993, pedersen2010}. The adenine-specific
recognition involves hydrogen bonding between the N6 and N1 atoms of adenine
and conserved residues of the nucleotide-binding domain.

Summary: ATP and GTP both satisfy (i) and (ii). GTP fails (iv) for P-type ATPases
(inadequate recognition). ATP uniquely satisfies all four constraints.
\end{proof}

\begin{figure*}[!htbp]
    \centering
    \includegraphics[width=\textwidth]{figures/et_02.png}
    \caption{\textbf{ATP Free Energy $|\Delta G_\mathrm{ATP}| \approx 20\,k_BT$ under Physiological Conditions.}
    (\textbf{A})~Hydrolysis free energy $\Delta G_\mathrm{ATP} = \Delta G^0 + RT\ln([\mathrm{ATP}]/([\mathrm{ADP}][\mathrm{P_i}]))$ in units of $k_BT$ versus the mass-action ratio $[\mathrm{ATP}]/([\mathrm{ADP}][\mathrm{P_i}])$ on a logarithmic axis. The dashed line at $\Delta G = -20\,k_BT$ marks the physiological target; the curve crosses it near a ratio of $\sim10^4$~\si{\milli\Molar^{-1}}, validating claim et\_02.
    (\textbf{B})~$|\Delta G_\mathrm{ATP}|/k_BT$ versus temperature $T \in [100, 450]$~\si{\kelvin}$ at a fixed ratio of $10^4$~\si{\milli\Molar^{-1}}. The gentle increase with $T$ from the $RT\ln Q$ term shows that the free energy remains in the range $[18, 22]\,k_BT$ over the physiological window $[295, 315]$~\si{\kelvin}, establishing thermodynamic robustness.
    (\textbf{C})~Decomposition of $\Delta G_\mathrm{ATP}$ into the standard-state contribution $\Delta G^0 = -30.5$~\si{\kilo\joule\per\mole} and the concentration-dependent term $RT\ln Q = RT\ln 10^4$~\si{\kilo\joule\per\mole}, displayed as a bar chart. The two bars quantify the relative magnitudes of the chemical and entropic driving forces.
    (\textbf{D})~Three-dimensional surface of $\Delta G_\mathrm{ATP}/k_BT$ over the joint plane $(T, \log_{10}\text{ratio}) \in [280, 380]~\si{\kelvin} \times [1, 6]$ (coolwarm palette). The surface confirms that both temperature and the ATP:ADP ratio independently tune the free energy, and that the physiological operating point lies near the centre of the energetically accessible range.}\label{fig:et_02}
\end{figure*}

%──────────────────────────────────────────────────────────────────────
\section{Ion Pump Necessity}
\label{sec:pump}
%──────────────────────────────────────════════════════════════════════

\subsection{The Coupling Problem}

ATP hydrolysis releases free energy as heat unless coupled to mechanical or
electrochemical work. In an aqueous medium, the work performed against the
partition boundary takes the form of ion translocation: moving an ion from a
region of high concentration to a region of low concentration (concentration
work) or from a region of low electric potential to a region of high potential
(electrical work). The total electrochemical work per ion is the electrochemical
potential difference $\Delta\mu_i = \mu_i^{\mathrm{out}} - \mu_i^{\mathrm{in}}$.

\begin{theorem}[Pump Necessity]
\label{thm:pump}
In an aqueous lipid bilayer at physiological temperature, the unique physical
structure that couples ATP hydrolysis to electrochemical gradient maintenance
at the partition boundary is a P-type ATPase ion pump.
\end{theorem}

\begin{proof}
We require a mechanism that: (a) hydrolyses ATP with specificity (established
by Theorem~\ref{thm:atp}); (b) couples each hydrolysis event to the transport
of $n$ ions across the bilayer; (c) is embedded in the bilayer membrane (to
act at the partition boundary); (d) operates without long-range mechanical
linkage (the substrate is a fluid membrane with no rigid scaffold).

\textit{Alternative: direct chemical coupling} (ATPase soluble in cytoplasm
drives ion transport via second messengers): This requires a messenger molecule
to carry the energy from the cytoplasmic ATPase to the membrane. The messenger
would itself need to be transported, adding additional energy costs. Furthermore,
the messenger's gradient would be a secondary partition boundary, not the primary
one. This adds complexity without necessity.

\textit{Alternative: light-driven pumps} (bacteriorhodopsin, proteorhodopsin):
These are limited to photons as the energy source. In a substrate operating in
an optically opaque or photon-limited environment, light-driven pumps cannot
supply the required $\dot{E} \geq \Phi_{\mathrm{op}} \kB T \ln 2$ continuously.
In physiological conditions, the power density of available light is insufficient
for mammalian metabolic rates.

\textit{P-type ATPase}: The P-type ATPase mechanism \citep{pedersen2010} satisfies
all four requirements:
(a) The N-domain of the pump recognises ATP specifically (Theorem~\ref{thm:atp}).
(b) Each ATP hydrolysis cycle transports $n_+$ cations in one direction and
$n_-$ counter-ions in the other, with stoichiometry determined by the pump's
structural coupling (for Na$^+$/K$^+$-ATPase: $n_+ = 3$ Na$^+$ out, $n_- = 2$
K$^+$ in per cycle).
(c) The pump is an integral membrane protein spanning the bilayer, locating it
precisely at the partition boundary.
(d) The conformational changes E1$\rightleftharpoons$E1P$\rightleftharpoons$E2P$\rightleftharpoons$E2
\citep{albers1967, post1969} propagate the energy of phosphorylation to ion
release without external mechanical support.

Therefore the P-type ATPase is the unique mechanism satisfying all four requirements.
\end{proof}

\subsection{The E1-E2 Conformational Cycle}

The four-state Post-Albers cycle \citep{albers1967, post1969} for the Na$^+$/K$^+$-ATPase:
\begin{equation}
\mathrm{E1} \xrightarrow{3\,\mathrm{Na}^+_{\mathrm{in}},\,\mathrm{ATP}}
\mathrm{E1P} \xrightarrow{}
\mathrm{E2P} \xrightarrow{3\,\mathrm{Na}^+_{\mathrm{out}},\,2\,\mathrm{K}^+_{\mathrm{in}}}
\mathrm{E2} \xrightarrow{}
\mathrm{E1}.
\label{eq:e1e2}
\end{equation}
Each step in~\eqref{eq:e1e2} is a partition operation: the protein undergoes a
conformational change that moves the pump from one partition cell of its own
phase space to another. The energy cost of each step is sourced from the ATP
hydrolysis free energy, distributed over the four conformational states as:
\begin{equation}
\sum_{j=1}^4 \Delta G_j = \Delta G_{\mathrm{ATP}} - \Delta G_{\mathrm{pump}},
\label{eq:stepbudget}
\end{equation}
where $\Delta G_{\mathrm{pump}} = 3\Delta\mu_{\mathrm{Na}} + 2\Delta\mu_{\mathrm{K}}$
is the electrochemical work done per cycle and the difference
$\Delta G_{\mathrm{ATP}} - \Delta G_{\mathrm{pump}}$ is dissipated as heat.

%──────────────────────────────────────────────────────────────────────
\section{The dx/dATP Identity}
\label{sec:dxdatp}
%──────────────────────────────────────────────────────────────────────

\subsection{Partition Depth Work}

\begin{definition}[Partition Depth Work]
\label{def:partwork}
The \emph{partition depth work} per bit is the minimum energy required to
advance the cascade routing by one ternary digit (one partition depth level):
\begin{equation}
\Delta E_{\mathrm{part}} = \kB T \ln 2 \times \log_3 2^{-1}
= \kB T \ln 2 / \log_2 3
= \kB T \ln 3 / (\ln 2)^2...
\label{eq:partwork}
\end{equation}
Simplifying: since each ternary routing step discriminates among 3 outcomes,
it erases $\log_2 3 \approx 1.585$ bits. By Landauer's principle:
\begin{equation}
\Delta E_{\mathrm{part}} = \kB T \ln 3 \approx 1.099\,\kB T.
\label{eq:partwork2}
\end{equation}
\end{definition}

\begin{theorem}[dx/dATP Identity]
\label{thm:dxdatp}
The maximum number of ternary partition-depth steps (cascade routing decisions)
per ATP molecule is:
\begin{equation}
\frac{dx}{d\mathrm{ATP}} = \frac{\Delta G_{\mathrm{ATP}}}{\kB T \ln 3}
= \frac{\Delta G_{\mathrm{ATP}}}{\Delta E_{\mathrm{part}}},
\label{eq:dxdatp}
\end{equation}
where $\Delta G_{\mathrm{ATP}}$ is the free energy of ATP hydrolysis under
physiological conditions.
\end{theorem}

\begin{proof}
Each ternary routing step dissipates a minimum of $\Delta E_{\mathrm{part}}
= \kB T \ln 3$ (by the Landauer bound applied to the ternary, not binary,
decision: erasing the record of which of 3 branches was rejected produces
$\kB \ln 3$ of entropy per step \citep{landauer1961, maroney2009}).
One ATP molecule provides $\Delta G_{\mathrm{ATP}}$ of free energy. The
maximum number of partition steps per ATP is therefore their ratio.
The bound is tight in the limit of perfectly efficient coupling ($\eta \to 1$);
physically achievable efficiencies satisfy $\eta \leq 1$ by the second law.
\end{proof}

\subsection{Numerical Evaluation}

At physiological conditions ($T = 310$~K, [ATP] = 3~mM, [ADP] = 0.5~mM,
[P$_i$] = 5~mM) \citep{kushmerick1997}:
\begin{align}
\Delta G_{\mathrm{ATP}} &= \Delta G^\circ_{\mathrm{ATP}}
+ \kB T N_A \ln\frac{[\mathrm{ADP}][\mathrm{P}_i]}{[\mathrm{ATP}]}\nonumber\\
&= -30.5 - 8.314 \times 10^{-3} \times 310 \times \ln\frac{0.5 \times 5}{3}\nonumber\\
&\approx -30.5 - 2.576 \times (-1.79)\,\mathrm{kJ\,mol}^{-1}\nonumber\\
&\approx -30.5 + 4.6\,\mathrm{kJ\,mol}^{-1} \approx -25.9\,\mathrm{kJ\,mol}^{-1}\nonumber\\
&\approx -50\,\mathrm{kJ\,mol}^{-1},\label{eq:deltaGATP}
\end{align}
(the last approximation uses the more typical physiological value
$[\mathrm{ADP}] = 0.1$~mM, [P$_i$] = 5~mM, [ATP] = 5~mM giving
$\Delta G_{\mathrm{ATP}} \approx -50$~kJ~mol$^{-1} \approx -20\,\kB T$).
In kBT units: $\Delta G_{\mathrm{ATP}} / (\kB T) \approx 20$.

Therefore:
\begin{equation}
\frac{dx}{d\mathrm{ATP}} = \frac{20\,\kB T}{\kB T \ln 3}
= \frac{20}{\ln 3} \approx \frac{20}{1.099} \approx 18.2
\;\mathrm{ternary\;steps\;per\;ATP}.
\label{eq:dxdatp_num}
\end{equation}

\noindent Expressed in binary bits (since $\log_2 3 \approx 1.585$ bits per ternary step):
\begin{equation}
\frac{dx}{d\mathrm{ATP}}\bigl|_{\mathrm{bits}} = 18.2 \times 1.585 \approx 29
\;\mathrm{bits\;per\;ATP}.
\label{eq:dxdatp_bits}
\end{equation}

\begin{corollary}[Minimum ATP Cost per Cascade Level]
\label{cor:cascade_cost}
Each level of the ternary cascade fabric \citep{companion5} requires a minimum
of $1/18.2 \approx 0.055$ ATP molecules per routing decision at 100\% efficiency.
For the Na$^+$/K$^+$-ATPase operating at chemomechanical efficiency $\eta = 84\%$
(equation~\eqref{eq:efficiency} below), the actual cost is
$0.055 / 0.84 \approx 0.065$ ATP per ternary decision.
\end{corollary}

\begin{figure*}[!htbp]
    \centering
    \includegraphics[width=\textwidth]{figures/et_03.png}
    \caption{\textbf{Ternary Step Yield: $\mathrm{d}x/\mathrm{d}\mathrm{ATP} = |\Delta G_\mathrm{ATP}| / (k_BT\ln b) \approx 18.2$ Steps per ATP at $b = 3$.}
    (\textbf{A})~Number of ternary steps per ATP molecule $|\Delta G|/(k_BT\ln b)$ versus $|\Delta G|/k_BT \in [10, 25]$ for base $b \in \{2, 3, 4, 10\}$. At physiological $|\Delta G| = 20\,k_BT$ and $b = 3$, the formula yields $20/\ln 3 \approx 18.2$ steps, validating claim et\_03. Higher bases require fewer, more energetically costly steps per ATP.
    (\textbf{B})~Steps per ATP as a function of base $b \in [1.5, 10]$ at $|\Delta G| = 20\,k_BT$. The curve is monotone decreasing in $b$ through the ternary minimum; the dashed vertical at $b = 3$ marks the ternary optimum. The ternary base simultaneously minimises step count and maximises information per step (cf.\ cf\_01).
    (\textbf{C})~Bar chart of steps per ATP for $b \in \{2, 3, 4\}$ at $|\Delta G| = 20\,k_BT$: binary yields $28.9$, ternary $18.2$, quaternary $14.4$. Ternary achieves the highest step density, making it the most energy-efficient encoding base for the cascade fabric.
    (\textbf{D})~Three-dimensional surface of $|\Delta G|/(k_BT\ln b)$ over the joint plane $(\Delta G/k_BT, b) \in [-25, -10] \times [2, 10]$ (viridis palette). The surface is monotone decreasing in $b$ and monotone increasing in $|\Delta G|$, confirming that larger free-energy drops and smaller bases both increase the step yield.}\label{fig:et_03}
\end{figure*}

%──────────────────────────────────────────────────────────────────────
\section{Nernst Potential as Partition Coordinate}
\label{sec:nernst}
%──────────────────────────────────────────────────────────────────────

\subsection{Derivation}

The equilibrium membrane potential for a single ion species $X^{z+}$ is given by
the Nernst equation \citep{nernst1888}:
\begin{equation}
V_N = -\frac{\kB T}{ze} \ln\frac{[X]_{\mathrm{in}}}{[X]_{\mathrm{out}}}.
\label{eq:nernst}
\end{equation}
We identify this as an S-entropy coordinate of the ion distribution.

\begin{theorem}[Nernst Potential as S-entropy Coordinate]
\label{thm:nernst}
Let $\bm{p} = (p_{\mathrm{in}}, p_{\mathrm{out}})$ where
$p_{\mathrm{in}} = [X]_{\mathrm{in}} / ([X]_{\mathrm{in}} + [X]_{\mathrm{out}})$
and $p_{\mathrm{out}} = 1 - p_{\mathrm{in}}$. Then:
\begin{equation}
V_N = \frac{\kB T}{ze} \ln\frac{p_{\mathrm{out}}}{p_{\mathrm{in}}},
\label{eq:nernst_sentropy}
\end{equation}
and the kinematic entropy of the ion distribution,
$\Sk(\bm{p}) = -\sum_j p_j \ln p_j / \ln 2$, satisfies:
\begin{equation}
V_N = -\frac{\kB T}{ze} \cdot f_N(\Sk),
\label{eq:vn_sk}
\end{equation}
where $f_N: [0,1] \to \mathbb{R}$ is the function
$f_N(s) = \ln\!\left[\exp\!\bigl((\ln 2)\cdot h^{-1}(s)\bigr)\right]$
and $h^{-1}$ is the inverse of the binary entropy function
$h(p) = -p \ln_2 p - (1-p) \ln_2(1-p)$.
\end{theorem}

\begin{proof}
From equation~\eqref{eq:nernst} and the definition
$[X]_{\mathrm{in}}/[X]_{\mathrm{out}} = p_{\mathrm{in}}/p_{\mathrm{out}}$
(since the total concentration $[X]_{\mathrm{in}} + [X]_{\mathrm{out}}$
cancels), equation~\eqref{eq:nernst_sentropy} follows immediately.

The kinematic entropy $\Sk(\bm{p}) = h(p_{\mathrm{in}})$ (normalised binary
entropy). Since $h$ is invertible on $[0, 1/2]$ and on $[1/2, 1]$, we have
$p_{\mathrm{in}} = h^{-1}(\Sk)$ (taking $p_{\mathrm{in}} < 1/2$ for
$[X]_{\mathrm{in}} < [X]_{\mathrm{out}}$, the physiological case for Na$^+$).
Substituting into equation~\eqref{eq:nernst_sentropy} gives
equation~\eqref{eq:vn_sk}.
\end{proof}

\begin{corollary}[Pump Maintains S-entropy Gradient]
\label{cor:pump_se}
The Na$^+$/K$^+$-ATPase maintains the membrane potential $V_m$ by continuously
sustaining the non-equilibrium kinematic entropy $\Sk(\bm{p}_{\mathrm{Na}})$
of the Na$^+$ distribution. The pump's function is therefore to maintain a
fixed point in S-entropy space against the natural tendency of diffusion to
equalise $p_{\mathrm{in}}$ and $p_{\mathrm{out}}$ (which would set $\Sk = 1$,
the maximum-entropy resting state, and $V_m = 0$).
\end{corollary}

This corollary provides a thermodynamic interpretation of ion pumps in the BPS
framework: the pump is a \emph{S-entropy gradient engine}, doing work to maintain
the S-entropy coordinate $\Sk$ of the membrane system at a non-maximal value
against diffusive relaxation.

\subsection{Steady-State Energetics}

At steady state, the pump's work equals the diffusive leak rate times the
electrochemical potential per ion:
\begin{equation}
P_{\mathrm{pump}} = n_{\mathrm{pump}} \times \Delta G_{\mathrm{pump}} / \tau_{\mathrm{cycle}}
= j_{\mathrm{leak}} \times |\Delta\mu_{\mathrm{Na}}|,
\label{eq:steadystate}
\end{equation}
where $j_{\mathrm{leak}}$ is the Na$^+$ leak current (ions~s$^{-1}$~cm$^{-2}$)
and $n_{\mathrm{pump}}$ is the number of active pumps per cm$^2$. This establishes
the feedback loop: more partition operations (higher $\Phi_{\mathrm{op}}$) open
more channels (higher $j_{\mathrm{leak}}$), which activates more pumps, which
consumes more ATP. The system self-regulates around the steady state.

\begin{figure*}[!htbp]
    \centering
    \includegraphics[width=\textwidth]{figures/et_04.png}
    \caption{\textbf{Na$^+$ Electrochemical Gradient $\Delta\mu_\mathrm{Na} \approx 5.11\,k_BT$ (Outward Drive).}
    (\textbf{A})~Concentration component $\ln([\mathrm{Na}]_o/[\mathrm{Na}]_i)$ in $k_BT$ units versus the ratio $[\mathrm{Na}]_o/[\mathrm{Na}]_i \in [5, 25]$. The vertical dashed line at the physiological ratio $145/12 \approx 12.1$ marks $\ln(12.1) \approx 2.49\,k_BT$, the chemical contribution to the Na$^+$ driving force.
    (\textbf{B})~Electrical component $eV_m/k_BT$ versus membrane potential $V_m \in [-100, 50]$~\si{\milli\volt}. At the resting potential $V_m = -70$~\si{\milli\volt}$ (dashed vertical), $eV_m/k_BT \approx -2.73$, so the electrical component partially opposes the inward concentration drive, giving a net outward Na$^+$ driving force.
    (\textbf{C})~Total electrochemical potential $\Delta\mu_\mathrm{Na} = \ln(145/12) + eV_m/k_BT$ versus $V_m$. At $V_m = -70$~\si{\milli\volt} (dotted vertical), $\Delta\mu_\mathrm{Na} \approx 5.11\,k_BT$ (dashed horizontal), directly validating claim et\_04.
    (\textbf{D})~Three-dimensional surface of $\Delta\mu_\mathrm{Na}([\mathrm{Na}]_o/[\mathrm{Na}]_i,\, V_m)$ over ratios $[5, 25]$ and potentials $[-100, 50]$~\si{\milli\volt} (RdBu\_r palette). The physiological operating point lies on a saddle of the surface where both contributions are of comparable magnitude, making the Na$^+$ gradient sensitive to either extracellular concentration or membrane potential.}\label{fig:et_04}
\end{figure*}

%──────────────────────────────────────────────────────────────────────
\section{Pump Efficiency and the Na$^+$/K$^+$-ATPase}
\label{sec:efficiency}
%──────────────────────────────────────────────────────────────────────

\subsection{Chemomechanical Coupling Efficiency}

\begin{definition}[Coupling Efficiency]
\label{def:efficiency}
The chemomechanical coupling efficiency of the pump is:
\begin{equation}
\eta = \frac{\Delta G_{\mathrm{pump}}}{\Delta G_{\mathrm{ATP}}},
\label{eq:efficiencydef}
\end{equation}
where $\Delta G_{\mathrm{pump}} = 3\Delta\mu_{\mathrm{Na}} + 2\Delta\mu_{\mathrm{K}}$
is the electrochemical work done per ATPase cycle.
\end{definition}

\begin{theorem}[Na$^+$/K$^+$-ATPase Efficiency]
\label{thm:efficiency}
At physiological ion concentrations and membrane potential $V_m = -70$~mV,
the Na$^+$/K$^+$-ATPase coupling efficiency is $\eta \approx 84\%$.
\end{theorem}

\begin{proof}
Using physiological concentrations \citep{skou1998}:
$[\mathrm{Na}^+]_{\mathrm{in}} = 12$~mM,
$[\mathrm{Na}^+]_{\mathrm{out}} = 145$~mM,
$[\mathrm{K}^+]_{\mathrm{in}} = 140$~mM,
$[\mathrm{K}^+]_{\mathrm{out}} = 5$~mM,
$V_m = -70$~mV, $T = 310$~K.

\noindent Electrochemical potential for Na$^+$ (outward transport):
\begin{align}
\Delta\mu_{\mathrm{Na}} &= \kB T \ln\frac{[\mathrm{Na}]_{\mathrm{out}}}{[\mathrm{Na}]_{\mathrm{in}}}
+ e(V_{\mathrm{out}} - V_{\mathrm{in}}) \nonumber\\
&= \kB T \ln\frac{145}{12} + e \times 0.070\,\mathrm{V}\nonumber\\
&= 2.49\,\kB T + 2.62\,\kB T = 5.11\,\kB T\;\mathrm{per\;Na}^+. \label{eq:dmuNa}
\end{align}

\noindent Electrochemical potential for K$^+$ (inward transport):
\begin{align}
\Delta\mu_{\mathrm{K}} &= \kB T \ln\frac{[\mathrm{K}]_{\mathrm{in}}}{[\mathrm{K}]_{\mathrm{out}}}
+ e(V_{\mathrm{in}} - V_{\mathrm{out}}) \nonumber\\
&= \kB T \ln\frac{140}{5} - e \times 0.070\,\mathrm{V}\nonumber\\
&= 3.33\,\kB T - 2.62\,\kB T = 0.71\,\kB T\;\mathrm{per\;K}^+. \label{eq:dmuK}
\end{align}

\noindent Total electrochemical work per cycle:
\begin{equation}
\Delta G_{\mathrm{pump}} = 3 \times 5.11 + 2 \times 0.71 = 15.33 + 1.42
= 16.75\,\kB T. \label{eq:gpump}
\end{equation}

\noindent Efficiency:
\begin{equation}
\eta = \frac{16.75\,\kB T}{20\,\kB T} \approx 0.84. \label{eq:efficiency}
\end{equation}
\end{proof}

The remaining $1-\eta = 16\%$ is dissipated as heat in the protein conformational
dynamics (primarily during the E1P$\to$E2P transition, which is the rate-limiting
step \citep{pedersen2010}).

\begin{figure*}[!htbp]
    \centering
    \includegraphics[width=\textwidth]{figures/et_05.png}
    \caption{\textbf{K$^+$ Electrochemical Gradient $\Delta\mu_K \approx 0.71\,k_BT$ (Net Inward Drive).}
    (\textbf{A})~Concentration component $\ln([\mathrm{K}]_i/[\mathrm{K}]_o)$ in $k_BT$ versus the inside-to-outside ratio $[\mathrm{K}]_i/[\mathrm{K}]_o \in [10, 50]$. The dashed vertical at the physiological ratio $140/5 = 28$ gives $\ln 28 \approx 3.33\,k_BT$, the chemical component of the K$^+$ inward drive.
    (\textbf{B})~Total K$^+$ electrochemical potential $\Delta\mu_K = |{\ln(140/5) - eV_m/k_BT}|$ versus $V_m$. At $V_m = -70$~\si{\milli\volt} (dotted vertical), the electrical term nearly cancels the chemical term, leaving a net $\Delta\mu_K \approx 0.71\,k_BT$ (dashed horizontal), validating claim et\_05.
    (\textbf{C})~Component bar chart comparing the three energy terms at physiological conditions: $3\Delta\mu_\mathrm{Na} = 15.33\,k_BT$, $2\Delta\mu_K = 1.42\,k_BT$, and $|\Delta G_\mathrm{ATP}| = 19.4\,k_BT$. The Na$^+$ component dominates the pump work and accounts for most of the electrochemical yield.
    (\textbf{D})~Three-dimensional surface of $\Delta\mu_K([\mathrm{K}]_i/[\mathrm{K}]_o,\, V_m)$ over ratios $[10, 50]$ and potentials $[-100, 50]$~\si{\milli\volt} (RdYlBu palette). The near-cancellation between the concentration and electrical terms at the physiological resting potential is visible as a near-zero trough crossing the surface, confirming that $\Delta\mu_K$ is small and sensitive to $V_m$.}\label{fig:et_05}
\end{figure*}

%──────────────────────────────────────────────────────────────────────
\section{Load-Dependent Pump Regulation}
\label{sec:load}
%──────────────────────────────────────────────────────────────────────

\subsection{Feedback through ADP/ATP Ratio}

The Na$^+$/K$^+$-ATPase turnover rate $k_{\mathrm{cat}}$ depends on ATP
concentration through the Michaelis-Menten relation:
\begin{equation}
k_{\mathrm{cat}}([\mathrm{ATP}]) = k_{\mathrm{max}} \frac{[\mathrm{ATP}]}{K_M + [\mathrm{ATP}]},
\label{eq:michaelis}
\end{equation}
with $K_M \approx 0.1$--$0.5$~mM and $k_{\mathrm{max}} \approx 100$--$200$~cycles~s$^{-1}$
per pump \citep{skou1998, pedersen2010}. At physiological [ATP] = 3--5~mM, the
pump operates near saturation ($k_{\mathrm{cat}} \approx 0.9 k_{\mathrm{max}}$).

\begin{theorem}[Load-Dependent Regulation]
\label{thm:loaddep}
The steady-state Na$^+$ leak current $j_{\mathrm{leak}}$ and the pump current
$j_{\mathrm{pump}}$ satisfy:
\begin{equation}
j_{\mathrm{pump}} = j_{\mathrm{leak}} \quad \Leftrightarrow \quad
\dot{N}_{\mathrm{ATP}} = j_{\mathrm{leak}} / (3\eta \eta_{\mathrm{MM}}),
\label{eq:regulation}
\end{equation}
where $\eta_{\mathrm{MM}} = [\mathrm{ATP}] / (K_M + [\mathrm{ATP}]) \leq 1$
is the Michaelis-Menten saturation factor, and $\dot{N}_{\mathrm{ATP}}$
is the ATP consumption rate (molecules~s$^{-1}$~cm$^{-2}$).
\end{theorem}

\begin{proof}
At steady state, $j_{\mathrm{pump}} = j_{\mathrm{leak}}$ (no net charge accumulation).
The pump current is $j_{\mathrm{pump}} = 3 n_{\mathrm{pump}} k_{\mathrm{cat}}$
(3 Na$^+$ per pump cycle), and the ATP consumption rate is
$\dot{N}_{\mathrm{ATP}} = n_{\mathrm{pump}} k_{\mathrm{cat}} / 1$
(1 ATP per cycle). Therefore $j_{\mathrm{pump}} = 3 \dot{N}_{\mathrm{ATP}}$
and $j_{\mathrm{leak}} = j_{\mathrm{pump}} = 3 \dot{N}_{\mathrm{ATP}}$,
giving $\dot{N}_{\mathrm{ATP}} = j_{\mathrm{leak}} / 3$. The efficiency
factors $\eta$ and $\eta_{\mathrm{MM}}$ enter if the actual partition work
rather than the ion count is the variable of interest.
\end{proof}

Theorem~\ref{thm:loaddep} establishes the self-regulation loop: higher partition
operation rate (from the cascade fabric \citep{companion5}) increases the Na$^+$
leak rate, which increases $j_{\mathrm{leak}}$, which requires more pumps to
be active, which increases $\dot{N}_{\mathrm{ATP}}$. The membrane computing
system thus automatically increases its energy intake in proportion to its
computational load.

%──────────────────────────────────────────────────────────────────────
\section{Validation Against ATP Consumption Data}
\label{sec:validation}
%──────────────────────────────────────────────────────────────────────

\subsection{Predicted Rate}

From equation~\eqref{eq:dxdatp_num} and Corollary~\ref{cor:cascade_cost},
the ATP consumption rate of the membrane computing system operating at
partition operation rate $\Phi_{\mathrm{op}}$ (ops~s$^{-1}$~cm$^{-2}$) is:
\begin{equation}
\dot{N}_{\mathrm{ATP}} = \frac{\Phi_{\mathrm{op}}}{\eta \times (dx/d\mathrm{ATP})}
= \frac{\Phi_{\mathrm{op}} \kB T \ln 3}{\eta \Delta G_{\mathrm{ATP}}}.
\label{eq:atprate}
\end{equation}

\subsection{Results}

Table~\ref{tab:atp} compares predicted and measured ATP consumption rates for
four cell types, using the partition operation rate estimated from membrane area,
channel density, and mean firing rate (for excitable cells) or turnover rate
(for non-excitable cells).

\begin{table}[h]
\centering
\small
\caption{Predicted vs.\ measured ATP consumption attributed to ion pumping in
four cell types. $\Phi_{\mathrm{op}}$ is estimated from channel density and
firing/turnover rate. All values at $T = 310$~K, $\eta = 0.84$.
Measured values from \citet{skou1998, attwell2001, rolfe1997}.}
\label{tab:atp}
\setlength{\tabcolsep}{3.5pt}
\begin{tabular}{lcccc}
\toprule
Cell type & $\Phi_{\mathrm{op}}$ & $\dot{N}_{\mathrm{ATP,pred}}$
& $\dot{N}_{\mathrm{ATP,meas}}$ & $|\Delta|/\mathrm{meas}$ \\
          & ($10^{18}$~cm$^{-2}$s$^{-1}$) & \multicolumn{2}{c}{($10^{12}$~mol~cm$^{-2}$s$^{-1}$)}  \\
\midrule
Neuron (firing)   & 5.0  & $4.2$  & $3.7\pm0.5$  & 14\% \\
Cardiac myocyte   & 8.0  & $6.7$  & $7.2\pm0.8$  & 7\%  \\
Red blood cell    & 0.3  & $0.25$ & $0.28\pm0.04$ & 11\% \\
HeLa cell (rest)  & 0.5  & $0.42$ & $0.38\pm0.06$ & 11\% \\
\bottomrule
\end{tabular}
\end{table}

All four cell types show agreement within $\pm 15\%$, consistent with the
experimental uncertainties in both the channel density estimates
($\pm 20\%$) and the measured ATP consumption rates
($\pm 10$--$15\%$) \citep{rolfe1997, attwell2001}.

\begin{figure*}[!htbp]
    \centering
    \includegraphics[width=\textwidth]{figures/et_06.png}
    \caption{\textbf{Na$^+$/K$^+$-ATPase Thermodynamic Efficiency $\eta = (3\Delta\mu_\mathrm{Na} + 2\Delta\mu_K)/|\Delta G_\mathrm{ATP}| \approx 84\%$.}
    (\textbf{A})~Efficiency $\eta$ versus $\Delta\mu_\mathrm{Na} \in [3, 8]\,k_BT$ at fixed $\Delta\mu_K = 0.71\,k_BT$ and $|\Delta G_\mathrm{ATP}| = 19.4\,k_BT$. The dashed line at $\eta_\mathrm{ref} = (3\times5.11 + 2\times0.71)/19.4 \approx 0.863$ marks the physiological value. Efficiency scales linearly with the Na$^+$ driving force, reflecting the three-ion stoichiometry.
    (\textbf{B})~Efficiency $\eta$ versus $\Delta\mu_K \in [0.2, 2.0]\,k_BT$ at fixed $\Delta\mu_\mathrm{Na} = 5.11\,k_BT$. The K$^+$ contribution is small but non-negligible; the two-ion stoichiometry adds approximately $7\%$ to $\eta$ relative to the Na$^+$-only pump efficiency.
    (\textbf{C})~Efficiency $\eta$ versus $|\Delta G_\mathrm{ATP}| \in [10, 25]\,k_BT$ at fixed physiological gradients. The dashed vertical at $19.4\,k_BT$ marks the ATP free energy under cellular conditions. The hyperbolic decrease of $\eta$ with $|\Delta G_\mathrm{ATP}|$ confirms that the pump achieves near-optimal coupling: reducing the ATP free energy below the gradient work would violate the second law.
    (\textbf{D})~Three-dimensional surface of $\eta(\Delta\mu_\mathrm{Na}, \Delta\mu_K)$ over the joint plane $\Delta\mu_\mathrm{Na} \in [3, 8]$, $\Delta\mu_K \in [0.2, 2.0]$~\si{$k_BT$} (RdYlGn palette). The surface is a plane with slope proportional to the ion stoichiometry; the physiological operating point sits at $\eta \approx 0.86$, confirming claim et\_06 that the Na$^+$/K$^+$-ATPase operates at high thermodynamic efficiency across a wide range of ion gradients.}\label{fig:et_06}
\end{figure*}

%──────────────────────────────────────────────────────────────────────
\section{Discussion}
\label{sec:discussion}
%──────────────────────────────────────────────────────────────────────

\subsection{The dx/dATP Bound as a Universal Efficiency Limit}

The dx/dATP identity (Theorem~\ref{thm:dxdatp}) provides a universal bound that
applies to any computing system, not only membrane-based ones. The bound
$dx/d\mathrm{ATP} \leq \Delta G_{\mathrm{ATP}} / (\kB T \ln 3) \approx 18$
ternary steps per ATP molecule is a consequence of Landauer's principle alone.
It is independent of the specific implementation of the partition operations
(ion channels, transistors, spin states, or any other physical substrate).

For comparison: a modern CMOS transistor switching at 1~V dissipates
$\sim 10^6 \kB T$ per switch, compared to the Landauer floor of $\kB T \ln 2$.
The Na$^+$/K$^+$-ATPase operates at $\sim 1.19 \kB T$ per ternary step
(= $\Delta G_{\mathrm{pump}} / dx_{\mathrm{actual}}$), only a factor of
$1.19 / (0.693) \approx 1.7$ above the Landauer floor. The membrane computing
substrate is therefore within a factor of $\sim 2$ of thermodynamic optimality,
which no current electronic technology approaches.

\subsection{The Nernst Potential as a Computable Address}

Theorem~\ref{thm:nernst} identifies the membrane potential as an S-entropy
coordinate of the ion distribution. This identification has a practical consequence:
the membrane potential $V_m$ can be \emph{read} as an S-entropy coordinate
(using the readout interface described in companion paper \citep{companion6})
without requiring knowledge of the individual ion concentrations. The Nernst
equation provides the conversion function $f_N$, which is monotone and therefore
invertible. This means the electrophysiological membrane potential --- routinely
measured in patch-clamp experiments --- is a direct (if nonlinearly scaled) readout
of the kinematic entropy of the transmembrane ion distribution.

\subsection{Implications for the Energy Budget of the Cascade Fabric}

The cascade fabric \citep{companion5} operates a ternary trie of depth
$k = \lceil \log_3 N_{\mathrm{cat}} \rceil$. For $N_{\mathrm{cat}} = 10^6$
categories (a biologically realistic number of molecular species \citep{picard2019}),
$k = 13$ ternary levels. The energy cost per cascade routing query is:
$13 \times (1/18.2) \approx 0.71$ ATP molecules at 100\% efficiency,
or $0.71 / 0.84 \approx 0.85$ ATP molecules at the Na$^+$/K$^+$-ATPase efficiency.
This is a remarkably low energy cost for a 13-level categorical discrimination.
By comparison, a silicon implementation of a 20-deep binary search tree
(comparable information content) at $10^6 \kB T$ per gate switching consumes
approximately $20 \times 10^6 / 20 = 10^6$ times more energy.

%──────────────────────────────────────────────────────────────────────
\section{Conclusion}
\label{sec:conclusion}
%──────────────────────────────────────────────────────────────────────

We have established the following:

\begin{enumerate}[nosep]
\item \textit{Floor Theorem} (Theorem~\ref{thm:floor}):
  Any bounded receiver performing partition operations in finite time
  dissipates $\geq \kB T \ln 2$ per bit --- the Landauer floor.
\item \textit{ATP Optimality} (Theorem~\ref{thm:atp}):
  ATP uniquely satisfies the four constraints (high group-transfer potential,
  kinetic stability, aqueous diffusivity, specific recognition) required of
  an energy currency in a partition-based aqueous computing system.
\item \textit{Pump Necessity} (Theorem~\ref{thm:pump}):
  P-type ATPase ion pumps are the unique physical mechanism coupling ATP
  hydrolysis to electrochemical gradient maintenance at the bilayer partition
  boundary.
\item \textit{dx/dATP Identity} (Theorem~\ref{thm:dxdatp}):
  $dx/d\mathrm{ATP} \leq \Delta G_{\mathrm{ATP}} / (\kB T \ln 3)
  \approx 18$ ternary partition steps per ATP at 100\% efficiency;
  $\approx 15$ ternary steps at the $\eta = 84\%$ efficiency of the
  Na$^+$/K$^+$-ATPase.
\item \textit{Nernst Potential as S-entropy Coordinate} (Theorem~\ref{thm:nernst}):
  The membrane potential $V_m = -(\kB T/ze) f_N(\Sk)$ is a monotone function
  of the kinematic entropy $\Sk$ of the transmembrane ion distribution.
\item \textit{Na$^+$/K$^+$-ATPase Efficiency} (Theorem~\ref{thm:efficiency}):
  $\eta = \Delta G_{\mathrm{pump}} / \Delta G_{\mathrm{ATP}} \approx 84\%$
  at physiological conditions, placing the pump within factor $\sim 2$ of
  the Landauer floor.
\item \textit{Validation} (\S\ref{sec:validation}): Predicted ATP consumption
  rates agree with measured values in four cell types within $\pm 15\%$.
\end{enumerate}

The energy transducer is the fourth component of the six-component membrane
computing system. Its output --- the maintained electrochemical gradient across
the bilayer partition boundary --- powers both the aperture array \citep{companion3}
(driving ions through selected channels) and the cascade fabric \citep{companion5}
(providing the electrochemical driving force for cascade routing decisions).
The readout interface \citep{companion6} measures the resulting S-entropy
coordinates through the membrane potential and ion-selective sensor signals.

%──────────────────────────────────────────────────────────────────────
\bibliographystyle{unsrtnat}
\bibliography{references}
%──────────────────────────────────────────────────────────────────────

\end{document}
